\documentclass{article}
\usepackage{amsmath}
\usepackage{amsfonts}
\usepackage{geometry}
\usepackage{enumitem}

\geometry{a4paper, margin=1in}

\begin{document}

\title{Practice Problems: Momentum and Impulse}
\author{}
\date{}
\maketitle

\section*{Problem 1: Graphical Analysis of Impulse}
A tennis player hits a $0.06 \text{ kg}$ tennis ball approaching her at $20 \text{ m/s}$. The force exerted by the racket on the ball over time is approximated by a trapezoidal Force vs. Time ($F-t$) graph:
\begin{itemize}
    \item The force increases uniformly from $0 \text{ N}$ to $600 \text{ N}$ over $0.002 \text{ s}$.
    \item The force remains constant at $600 \text{ N}$ for the next $0.004 \text{ s}$.
    \item The force decreases uniformly from $600 \text{ N}$ to $0 \text{ N}$ over the final $0.002 \text{ s}$.
\end{itemize}
\begin{enumerate}[label=(\alph*)]
    \item Calculate the magnitude of the impulse delivered to the ball by the racket.
    \item Assuming the initial velocity of the ball was in the negative direction, what is the final velocity of the ball after the hit?
\end{enumerate}

\newpage

\section*{Problem 2: Bouncing vs. Sticking (Advanced Concept)}
A $0.2 \text{ kg}$ snowball is thrown horizontally at $15 \text{ m/s}$ toward a stationary $1.8 \text{ kg}$ dummy target resting on a frictionless ice surface. Consider two different scenarios:
\begin{itemize}
    \item \textbf{Scenario A:} The snowball is loose and completely sticks to the target upon impact.
    \item \textbf{Scenario B:} The snowball is hard ice and rebounds horizontally back in the opposite direction at $5 \text{ m/s}$.
\end{itemize}
\begin{enumerate}[label=(\alph*)]
    \item Without performing specific calculations, use the Impulse-Momentum Theorem to explain in which scenario the dummy target will acquire a greater final speed.
    \item Calculate the final speed of the target in both scenarios to verify your conclusion.
\end{enumerate}

\newpage

\section*{Problem 3: Algebraic Derivation}
Two skaters are initially standing at rest on a frictionless ice surface. Skater A has mass $M$ and Skater B has mass $2M$. Skater A pushes Skater B, causing Skater B to slide to the right with a final speed $v$.
\begin{enumerate}[label=(\alph*)]
    \item Derive an expression for the final velocity of Skater A (including direction) in terms of $M$ and $v$.
    \item During the push, how does the magnitude of the impulse experienced by Skater A compare to the magnitude of the impulse experienced by Skater B? Explain your answer using fundamental physics principles.
\end{enumerate}


\newpage

\section*{Problem 4: System Boundaries (Conceptual)}
A $30 \text{ kg}$ child stands on a $10 \text{ kg}$ sled, both initially at rest on a frictionless frozen lake. The child throws a $2.0 \text{ kg}$ block of ice horizontally forward with a speed of $10 \text{ m/s}$ relative to the ground. 

For each of the following defined systems, state whether the momentum of the system is conserved during the throw and briefly explain why:
\begin{enumerate}[label=(\alph*)]
    \item \textbf{System 1:} Only the thrown block of ice.
    \item \textbf{System 2:} The child and the sled (not including the block of ice).
    \item \textbf{System 3:} The child, the sled, and the block of ice.
\end{enumerate}
\newpage

\section*{Problem 5: Multi-stage Process}
A $4.95 \text{ kg}$ wooden block rests on a horizontal floor with friction. A $0.05 \text{ kg}$ bullet traveling horizontally at $400 \text{ m/s}$ strikes the block and instantly embeds itself within it (the collision time is negligible). After the collision, the block (with the bullet inside) slides across the floor for $2.0 \text{ s}$ before coming to a stop.
\begin{enumerate}[label=(\alph*)]
    \item Why is it valid to assume momentum is conserved for the "bullet + block" system during the extremely short collision, despite the presence of friction from the floor? Calculate the initial sliding speed of the block immediately after the collision.
    \item During the sliding phase (from just after the collision until it stops), what is the total change in momentum ($\Delta \vec{p}$) of the "bullet + block" system? What external force provides the impulse that causes this change, and what is the magnitude of this force?
\end{enumerate}
\vspace{4cm}

\end{document}
